% =========================================================
% TRACKED-CHANGES copy. The 2026-05-03 Garg-review revision
% is highlighted in red:
%   \new{...}     wraps inserted text
%   \removed{...} wraps deleted text (red strikethrough)
% Comment markers above each red region indicate which Garg
% comment they address (C2/C3/C4/C5).
% =========================================================

\paragraph{Model merging is everywhere, but nobody knows the
limit.} Low-Rank Adaptation~\cite{hu2022lora} has made it routine
to maintain dozens of task-specific fine-tunes per base model,
and it has made \emph{merging} those fine-tunes into a single
deployable model the dominant practical question. A now-standard
workflow: fine-tune $T$ LoRA adapters on $T$ downstream tasks,
combine them post-hoc, serve one model. Task
Arithmetic~\cite{ilharco2023task},
TIES~\cite{yadav2023ties}, DARE~\cite{yu2024dare},
KnOTS~\cite{stoica2025knots},
DO-Merging~\cite{domerging2025},
Core~Space~\cite{panariello2025core},
Task-Vector Quantization~\cite{kim2025tvq},
and 1-bit-Merging~\cite{bit1merging2025} all compete on this
problem. Each paper proposes an encoder, tunes some
hyper-parameters, and reports benchmark numbers. None answers
the prior question: given a storage budget of $R$ bits per
merged model, what is the \emph{best possible} merging
algorithm, and how far are existing methods from it?

% C2: the two new motivation sentences before the formalism are
% wrapped in \new{...}; the deleted sentence "Given such a tuple,"
% and "The figure of merit is the worst-task distortion ... only
% as good as its worst task." are shown via \removed{...}.
\paragraph{The rate-distortion question.} We formalize this as
multi-source lossy source coding. \new{The first design choice is
the figure of merit. A merged model that scores well on most
downstream tasks but collapses on one is not a useful
deployment target — quality should be measured by the
\emph{worst} task, not the average. Once this choice is fixed,
the formalism is direct.} Fix $T$ tasks, ambient dimension $d$,
and LoRA rank $r\leq d$. Each task is characterized by a pair
$(\tau_t, H_t)$: a task vector $\tau_t\in\R^d$ (the LoRA
update) and a positive-semidefinite $H_t$ of rank $r$ (the
task's Hessian at $\tau_t$, or a Fisher surrogate) with
$\tau_t^\top H_t\tau_t\leq B^2$. \removed{Given such a tuple, an}\new{An}
encoder maps $(\tau_1,\ldots,\tau_T)$ into $R$ bits; a decoder
reconstructs a single merged model $w^\star\in\R^d$. Per-task
distortion is $D_t(w) := (w-\tau_t)^\top H_t (w-\tau_t)$, the
second-order loss inflation from using $w$ in place of $\tau_t$.
\removed{The figure of merit is the \emph{worst-task} distortion
$\max_{t\leq T} D_t(w^\star)$ — a merged model is only as good
as its worst task.} The rate-distortion function
$\mathcal{D}^\star(T,d,r,B,R)$ is the infimum over encoders and
decoders of expected max-distortion \new{$\max_{t\leq T} D_t(w^\star)$}
at rate $R$.

Two features separate this problem from classical lossy
source coding. First, the encoder observes $T$ sources rather
than one; multi-description coding~\cite{elgamalcover1982md}
addresses related questions but under \emph{average} distortion,
which Lemma~\ref{lem:id} (below) shows is strictly weaker.
Second, the low-rank structure of $H_t$ means distortions are
measured in task-specific rank-$r$ subspaces; the effective
dimension is the combinatorial quantity $\deff := \E[\rank
(\sum_t P_{V_t})]$, which satisfies $\deff \leq \min(Tr, d)$,
with the bound attained almost surely when the $V_t$'s are
drawn iid from the Stiefel manifold.

\paragraph{Contributions.} We give the first rate-distortion
theorem for model merging, matching lower and upper bounds up
to universal constants in the generic regime.

% C3: each of the 5 contribution bullet titles was rewritten to
% lead with a plain-English claim. Wrap the new bold titles
% (and any newly-added prose tail) in \new{...}.
\begin{enumerate}[leftmargin=*,topsep=2pt,itemsep=2pt]
\item \textbf{\new{An algebraic identity reducing max-distortion
merging to average-distortion quantization plus an irreducible
floor} (Lemma~\ref{lem:id}).} For every admissible $(\tau,H)$
and every $w\in\R^d$,
\[
\max_t D_t(w)\;\geq\; (w-\tauH)^\top\bar H(w-\tauH)
+ \tfrac{1}{T}\sum_t (\tau_t-\tauH)^\top H_t(\tau_t-\tauH),
\]
where $\tauH = \bar H^+\cdot T^{-1}\sum_t H_t\tau_t$ is the
$H_t$-weighted centroid and $\bar H = T^{-1}\sum_t H_t$. The
identity is the technical bridge from Shannon's single-source
rate-distortion machinery to the multi-source max-distortion
setting\new{: the first term is what bits buy, the second is what
no rate can erase}.

\item \textbf{\new{A closed-form expression for the irreducible
floor in terms of subspace overlap} (Lemma~\ref{lem:floor-general}).}
Under a Stiefel-random hard distribution $P^\star$ on
$(\tau, H)$, the floor reduces to $B^2(1 - \deff/(Tr))$
exactly. The interpolation between $T=1$ (floor $0$, all bits
go to compression) and $V_t$ mutually orthogonal (floor $0$
again, reconstructive task arithmetic is lossless at infinite
rate) is explicit\new{; positive overlap raises the floor}.

\item \textbf{\new{A rate-distortion lower bound on every merging
algorithm} (Theorem~\ref{thm:general}).} Under $P^\star$ with
arbitrary positive-definite spectra $(D_t)_{t=1}^T$, \new{no
encoder/decoder pair at rate $R$ bits can do better than}
\begin{equation*}
\mathcal{D}^\star(T,d,r,B,R)\;\geq\;
B^2\!\left(1 - \tfrac{\deff}{Tr}\right) +
c_{\mathrm{TQ}}\cdot\tfrac{B^2\deff}{Tr}\cdot 2^{-2R/\deff}.
\end{equation*}
The proof composes Yao's minimax principle, Lemma~\ref{lem:id},
a reduction to the $\range(\bar H)$ subspace, and an
$O(d)$-equivariance argument that works regardless of
task-dependent spectra.

\item \textbf{\new{An explicit algorithm that matches the lower
bound up to a universal constant} (Theorem~\ref{thm:achv}).}
In the floor-zero Stiefel regime ($Tr\leq d$, $V_t$ drawn iid
from the Stiefel manifold), \new{randomized orthogonal mixing of
the $\bar H^{1/2}$-scaled centroid followed by uniform scalar
quantization achieves}
$\E[\max_t D_t(w^\star)]\leq C\cdot B^2\cdot 2^{-2R/(Tr)}$
for a universal constant $C = O(T)$. Combined with
Theorem~\ref{thm:general},
$\mathcal{D}^\star = \Theta\!\left(B^2\cdot 2^{-2R/(Tr)}\right)$
in this regime. The construction reduces to
TurboQuant~\cite{zandieh2025turboquant} at $T=1$.

\item \textbf{Numerical validation \new{of the predicted slopes
and constants}.} Monte Carlo experiments across $T\in\{2,3,4,8\}$,
$d\in\{128,512,2048\}$, $r\in\{4,8,16\}$, with both isotropic
and heavily anisotropic task-dependent spectra, confirm the
predicted slope $-2R/\deff$ on $\log_2\E[\text{excess}]$ to
within bootstrap 95\% CI $\pm 0.01$ across 1000 trials per
configuration. The achievability constant $C$ is empirically
$\approx 11$--$13$, matching the analytical bound $Tc^2/3$.
\end{enumerate}

\paragraph{What this means for practice.} The bound gives
practitioners a calibration curve. Any merging algorithm —
existing or future — can be graded by its gap to
$\mathcal{D}^\star$. Two qualitative implications are already
visible from the theorem statement. (i) The floor
$B^2(1 - \deff/(Tr))$ is irreducible: when $V_t$'s overlap
substantially ($\deff\ll Tr$), no amount of rate closes it.
This recovers and explains the recent large-scale empirical
finding of~\cite{systematic2025merging} that subspace-based
merging methods fail at LLM scale — random LoRA checkpoints
have substantial subspace overlap, so the floor is bounded
away from zero. (ii) Additional rate helps at the
$2^{-2R/\deff}$ rate, matching what TurboQuant achieves in
the single-source case. A practitioner storing a merged
model at 4 bits per coordinate operates roughly
$4\deff$-bits-below the infinite-rate floor; doubling to
8 bits gains an $\approx 16\times$ reduction in compression
excess.

% C4: the entire worked-example paragraph below is new.
\new{\emph{A worked example.} Consider a Llama-3-8B base model with
$T=4$ task-specific LoRAs at rank $r=16$. The per-layer
effective dimension satisfies $\deff\leq Tr = 64$. When the
four task subspaces $V_t$ are mutually orthogonal, the floor
vanishes and quality is rate-limited; when they fully overlap
($\deff = 16$), the floor jumps to $\approx 0.75 B^2$, a
ceiling no merging algorithm can break through. The two
regimes call for opposite operational responses: rate-limited
deployment benefits from spending the bit budget on
quantization precision, while overlap-limited deployment
benefits from serving tasks from separate adapters rather than
merging at all. The theorem identifies the crossover.}

\paragraph{Scope and limitations.} The matching upper bound
in Theorem~\ref{thm:achv} is proved under the generic
Stiefel-random hard distribution $P^\star$ with $Tr\leq d$.
In the floor-positive regime (shared or heavily overlapping
$V_t$), an explicit linear encoder achieves an excess rate of
$2^{-2R/(\deff+|A|-1)}$ (where $|A|$ is the size of the
active set at the Chebyshev center), and we verify this
numerically for $T\in\{2,3,4\}$; closing the exponent gap to
the lower bound's $2^{-2R/\deff}$ appears to require
non-linear / random-coding encoders and is left open
(Remark~\ref{rem:sharedv-gap}). The bound is stated in the
MSE / quadratic-loss setting; the extension to cross-entropy
LLM loss via local Fisher-quadratic approximation is
discussed in~\S\ref{sec:discussion}.

% C5: the phrase " in six steps" was removed per Garg's flag.
\paragraph{Paper organization.}
\S\ref{sec:setup} introduces the formal rate-distortion model
for merging and the admissible hard distributions.
\S\ref{sec:lower} proves the lower bound
(Theorem~\ref{thm:general})\removed{ in six steps}, with the
max-$\geq$-avg identity and closed-form floor as the central
tools. \S\ref{sec:achv} gives the matching achievability
construction and analyses. \S\ref{sec:exp} reports synthetic
experiments on rank-$r$ LoRA-like ensembles validating the
theoretical slopes and constants. \S\ref{sec:discussion}
discusses extensions to cross-entropy loss, the open
shared-$V$ exponent gap, and implications for practical
merging algorithm design.
